\begin{abstract}
Language model agents combine learned representations with external actions.
We study how memory and feedback affect planning across interactive environments.
Our analysis separates the policy that proposes an action from the mechanism
that records observations. This separation supports reproducible comparisons
under a fixed interaction budget.
\end{abstract}

\section{Introduction}
Interactive systems must choose useful actions while maintaining a compact
representation of previous observations. A successful action can provide both
immediate progress and information about the environment. Prior research on
language models, reinforcement learning, and software tools motivates a common
evaluation protocol. We consider tasks that require several dependent actions,
where the outcome of one step changes the information available to the next.
The literature provides complementary perspectives on this problem: CITATIONS

\section{Method}
We represent the environment as a sequence of observations and actions.
At each step the agent receives an observation, retrieves relevant memory,
and proposes an action using a fixed policy. The memory contains a summary
of earlier interactions and a record of outcomes. The policy receives the
same total context budget in every condition. We compare a chronological
memory with a retrieval-based memory that selects observations according
to their relevance to the current task.

Feedback is recorded after each action. Successful actions contribute to
progress, while unsuccessful actions provide evidence about invalid plans.
We separate feedback available during interaction from the final evaluation.
This prevents a policy from receiving information that would be unavailable
in normal use. Each run starts from the same environment state and uses
an independent interaction history.

\section{Evaluation}
We measure completion, interaction cost, and consistency across repeated
runs. Completion measures whether the agent reaches the requested state.
Interaction cost counts actions and observations under a common budget.
Consistency captures variation across independent runs with equivalent
initial conditions. We report all three quantities because a higher
completion rate may require substantially more interaction.

The protocol includes tasks with delayed feedback and tasks where several
actions have similar immediate outcomes. These conditions help distinguish
short-term response quality from sustained planning. We inspect action
sequences to understand when memory helps recover from an unsuccessful step.
The analysis considers both aggregate outcomes and the distribution of
interaction lengths. Future experiments can extend the protocol to richer
environments while preserving the separation between feedback and evaluation.
